%! Displaying and Analyzing Categorical Data — Unit 3, Lesson 3.1 — Student Packet Cover Sheet
% THREE packet rows — Warm-Up, Guided Notes & Practice, Homework. The homework is the individual
% practice, assigned at Close & Assign and done outside class, and it is scored 1 to 5. The cover
% is a contents list, with no Score column.
\documentclass[12pt]{article}
\usepackage{saar-article}
\usepackage{saar-boxes}
\usepackage{ltablex}
\keepXColumns
% Measured banner: the band grows to fit the title block, so a title long enough to wrap grows
% the band rather than running out of it (the stats course's \coverbanner; saar-article does not
% carry one, so it is defined here — every cover copies this block verbatim).
\newsavebox{\bannerbox}
\newlength{\bannerht}
\newlength{\coverbannerpad}
\setlength{\coverbannerpad}{0.16in}       % breathing room above and below the text
\newcommand{\coverbanner}[2]{%
  \sbox{\bannerbox}{%
    \begin{minipage}{\dimexpr\paperwidth-1.4in\relax}%
      \centering
      {\color{white}\textbf{\LARGE \CourseName}}\\[4pt]
      {\color{frostmid}\large\textbf{#1}}\\[6pt]
      {\color{white}\textbf{\large #2}}%
    \end{minipage}}%
  \setlength{\bannerht}{\dimexpr\ht\bannerbox+\dp\bannerbox+2\coverbannerpad\relax}%
  \noindent\begin{tikzpicture}[remember picture, overlay]
    \fill[cerulean] (current page.north west)
      rectangle ([yshift=-\bannerht]current page.north east);
    \node[anchor=north, inner sep=0pt]
      at ([yshift=-\coverbannerpad]current page.north) {\usebox{\bannerbox}};
  \end{tikzpicture}\par
  % The text body starts 0.6in below the page edge (geometry top=0.6in), so reserve only what
  % the band overruns that, plus a gap under the band.
  \vspace*{\dimexpr\bannerht-0.6in+0.16in\relax}%
}

\begin{document}

% Full-bleed cerulean banner, sized to the title block.
\coverbanner{Unit 3}{Lesson 3.1 \quad Displaying and Analyzing Categorical Data}

\vspace{0.06in}
% NAMESTRIP: this is the ONLY component that carries the name row.
\namedateperiod
\vspace{0.18in}

% GRADUAL RELEASE: the vocabulary is taught in the guided notes, so the targets name it
% plainly. The standard codes (PS.DS.4a-d) stay in the teacher-facing plan.
\begin{learningtargetbox}
\small
\begin{enumerate}[leftmargin=1.5em, itemsep=5pt]
  \item \ldots{} read a bar graph of one categorical variable in context, as a frequency or a
        \textbf{relative frequency}.
  \item \ldots{} compare groups with \textbf{side-by-side} and \textbf{segmented} bar graphs.
  \item \ldots{} build a \textbf{two-way table} from a summary and read its joint counts and
        \textbf{marginal frequencies}.
  \item \ldots{} calculate and interpret marginal and \textbf{conditional relative frequencies},
        and name the whole each percent is taken of.
  \item \ldots{} compare groups of different sizes with conditional percents, not raw counts ---
        and explain why.
\end{enumerate}
\end{learningtargetbox}

\vspace{0.14in}

\begin{tocbox}
\small
\renewcommand{\arraystretch}{1.5}
\begin{tabularx}{\linewidth}{@{} c l X @{}}
\rowcolor{frost}
\textbf{\#} & \textbf{Component} & \textbf{Description} \\
\midrule
1 & Warm-Up & The start of Unit~3: percents of a number, fractions as percents, and reading a
             simple table \\
2 & Guided Notes \& Practice & How Oakdale High students get to school: bar graphs, two-way
             tables, and conditional percents, with a vote on a newsletter's claim --- then a
             library's streaming survey we work together \\
% Row 3 is the lesson's GRADED practice, assigned at the close. Packet pages are the default;
% the form may change on a given day, decided at Close & Assign. The row never changes.
3 & Homework & Pet adoptions at a shelter: a two-way table, the percent of which whole, and a
             claim built on counts. \textbf{Scored 1 to 5, assigned today, due the first class
             after two study halls.} \\
\end{tabularx}
\end{tocbox}

\vspace{0.14in}

% Keep in Mind carries the lesson's CONTENT — the rule, the test, the trap — never its process.
\begin{remindbox}
\small
A \textbf{relative frequency} is a count divided by a total. In a \textbf{two-way table}, a
\textbf{marginal} relative frequency divides a Total by the whole table; a \textbf{conditional}
relative frequency divides an inside cell by the total of one group. \textbf{The test:} in ``Of
the [group], what percent \ldots?'' the words after \emph{of} name the whole, and that total is
the denominator. \textbf{Watch out:} a larger count does not mean more likely. To compare groups
of different sizes, compare conditional percents, never raw counts.
\end{remindbox}

\end{document}
